\documentclass[conference]{IEEEtran}
\IEEEoverridecommandlockouts
\usepackage[utf8]{inputenc}
\usepackage{parskip}
\usepackage{cite}
\usepackage{amsmath,amssymb,amsfonts}
\usepackage{algorithmic}
\usepackage{graphicx}
\usepackage{textcomp}
\usepackage[most]{tcolorbox}
\usepackage{xcolor}
\usepackage{url}
\usepackage{placeins}
\usepackage{tabularx}
\usepackage{makecell}
\usepackage{caption}

\title{Guardrailed LLM Agent Orchestration for Network Change Automation: A Comparative Implementation Study of n8n, Langflow and LangGraph}

\author{
\IEEEauthorblockN{Sandra Ilievska}
\IEEEauthorblockA{
\textit{Faculty of Computer Science and Engineering} \\
\textit{Ss. Cyril and Methodius University in Skopje} \\
Skopje, North Macedonia \\
sandra.ilievska@students.finki.ukim.mk}
\and
\IEEEauthorblockN{Sonja Filiposka}
\IEEEauthorblockA{
\textit{Faculty of Computer Science and Engineering} \\
\textit{Ss. Cyril and Methodius University in Skopje} \\
Skopje, North Macedonia \\
sonja.filiposka@finki.ukim.mk}
\and
\IEEEauthorblockN{Roman Łapacz}
\IEEEauthorblockA{
\textit{Poznan Supercomputing and Networking Center - PSNC} \\
Poznan, Poland \\
romradz@man.poznan.pl}
}
% TODO(Sandra): insert Roman Lapacz's correct affiliation email in the block above before submission

\begin{document}

\maketitle

\begin{abstract}
Production networks tolerate close to zero error, yet LLM agents are probabilistic by construction. This tension, not the choice of tooling, is the central obstacle to network-change automation: operators trust an agentic proposal only when a human stays in the loop and every step is bounded by a deterministic guardrail that inspects its structured output before anyone downstream sees it. We implement an identical ten-step, guardrail-checked firewall-provisioning pipeline, from intent parsing through conditional rollback, in three orchestration platforms: n8n, Langflow and LangGraph, each LLM-driven step returning structured JSON that is independently re-checked for schema validity and, where available, against ground truth computed outside the model. We run four fixed intents ten times each on LangGraph and n8n, and ten times against Langflow's connected human-approval-to-rollback segment, whose remaining steps could not be wired end-to-end in time. The results surface two failure modes: LLM stochasticity, which rejects deliberately safe requests at a non-trivial rate, and a platform-level architectural gap, where n8n's HTTP wiring silently drops a check LangGraph's code-first state model enforces by construction, letting a policy-violating request reach apply in 8 of 10 trials. A second, larger reasoning model (GLM-5.1) shows none of the false-rejection bias on the same harness, indicating it is a property of the smaller local model, not the guardrail design. Code, guardrail logic and raw experiment logs are publicly available.
\end{abstract}

\begin{IEEEkeywords}
LLM agents, guardrails, agentic orchestration, network automation, human-in-the-loop, n8n, Langflow, LangGraph
\end{IEEEkeywords}

\section{Introduction}

A firewall rule applied to the wrong subnet, or a rollback that silently fails to trigger, has consequences that a malformed chatbot reply does not. This is the property that separates network change automation from most other LLM agent applications, and it sets the requirement this paper takes as its starting point: an agentic pipeline for network operations must assume the model will occasionally be wrong, and must be engineered so that a wrong output is caught before it propagates, not after.

Two design choices follow from this requirement. First, a human must approve the change before it is applied, because no current guardrail can substitute for operator judgment on business risk. Second, and less commonly discussed, every intermediate agent in the pipeline must be constrained by a guardrail of its own. Without this, an error introduced by the first agent -- for instance, a misread source subnet -- propagates silently through validation, policy checking and planning, and only surfaces to the human at the approval step, by which point the reasoning chain that produced it is opaque. We therefore treat per-agent guardrails, not the choice of orchestration engine, as the primary independent variable of this study.

Whether these guardrails live inside a hand-written agent or are assembled from nodes in a general-purpose orchestration platform matters less than is often assumed. Both paths cost roughly the same engineering effort, spent on the same two things: a strict output contract for each agent, and a deterministic check that enforces it. What does differ across platforms is how naturally that contract can be expressed, and how easily a check can be silently omitted. This paper investigates that question empirically. We implement one fixed use case -- natural-language-driven firewall-rule provisioning, using NetBox as inventory and a mock firewall API -- identically in n8n, Langflow and LangGraph, and instrument every agent step with the same guardrail logic. We then run repeated trials of four fixed intents against the fully-wired LangGraph and n8n pipelines, and repeated trials of a fixed approval request against Langflow's connected human-approval-to-rollback segment (Section V explains why Langflow's scope is narrower), to measure both the failure rate and, where it differs across platforms, its cause.

The remainder of this paper is organized as follows. Section II reviews related work on LLM agents in network management. Section III describes the three platforms. Section IV presents the use case and the guardrail architecture shared across implementations. Section V describes the experimental methodology. Section VI reports the results. Section VII discusses the two classes of failure observed and their implications. Section VIII concludes.

\section{Related Work}

Recent surveys document the growing use of LLMs in network and service management, spanning configuration, monitoring and troubleshooting \cite{hong2025survey}, communication/mobile/edge domains \cite{boateng2025survey}, and coordination topologies for agentic systems in general \cite{futureinternet2026}.

Closer to our use case, purpose-built multi-agent systems apply LLMs to network management directly. Confucius models network workflows as directed acyclic graphs with retrieval-augmented generation and explicit human/model interaction \cite{wang2025confucius}. Brodimas \textit{et al.} combine agentic AI with intent-based networking to translate natural-language intents into infrastructure actions \cite{brodimas2025intent}. Zaid \textit{et al.} report a 96.7\% success rate translating user intent into optical-network operations without fine-tuning \cite{zaid2025multiagent}.

None of these systems reports trial-to-trial guardrail stability under repeated, identical input -- the central measurement of this paper. To our knowledge, no study benchmarks general-purpose orchestration platforms against one another on an identical, safety-critical automation pipeline, or isolates whether a guardrail failure originates in model stochasticity or in the orchestration layer's wiring. This is the gap this paper addresses.

\section{Orchestration Platforms Under Study}

\subsection{n8n}
n8n is a visual, node-based workflow automation platform. An LLM step is a chain or agent node wired to arbitrary downstream nodes -- HTTP requests, conditional (If) nodes, code nodes. A guardrail is a separate HTTP call to a check endpoint followed by an If node branching on its response; as Section VII shows, this also makes it possible to wire the If node against the wrong field without any static check catching the mistake. n8n's Wait node is a documented HITL primitive: it pauses execution and exposes a resumable form URL that an operator (or, here, an automated approval) posts back to.

\subsection{Langflow}
Langflow is an open-source, Python-based visual framework, wired on a canvas like n8n, but each component is a thinner wrapper around Python/LangChain primitives. HITL support is comparatively immature: it is exposed only through a background-execution v2 API (\texttt{/api/v2/workflows}) that suspends a run at a named component and resumes it via a separate endpoint. The more commonly documented synchronous build endpoint (\texttt{/api/v1/build}) does not honor this suspension at all and silently executes both branches of a pending decision -- a pitfall documented in Section V.

\subsection{LangGraph}
LangGraph is a code-first runtime in which a workflow is an explicit state graph: nodes are Python functions on a shared, typed state object, and edges (including conditional edges) are ordinary control flow. HITL is a first-class primitive (\texttt{interrupt}), and, critically here, a guardrail check is a Python function call whose result and the upstream agent's output are both directly addressable fields of the same state object -- removing the wiring step that introduces risk in n8n and Langflow, at the cost of requiring code rather than a canvas.

\begin{table}[!t]
\caption{Qualitative Comparison Criteria}
\label{tab:criteria}
\centering
\footnotesize
\begin{tabularx}{\columnwidth}{@{}l X@{}}
\hline
\textbf{Criterion} & \textbf{Observation} \\
\hline
Workflow definition & Canvas (n8n, Langflow) vs.\ code (LangGraph): faster to start vs.\ easier to diff/review \\
Guardrail wiring & LangGraph co-locates output and decision field; canvas platforms split them across nodes (VII.A's gap) \\
HITL maturity & n8n Wait: stable, documented. Langflow v2 API: works, silent-bypass pitfall (III.B). LangGraph \texttt{interrupt}: first-class \\
Pre-/post-change checks & Shared Python module on all three: LangGraph imports directly, n8n/Langflow via HTTP wrapper \\
Rollback mechanism & Implemented identically (Table~\ref{tab:usecase}, step 10) on all three \\
Observability & First-class state field in LangGraph; n8n/Langflow need post-hoc per-node reads \\
Implementation effort & Comparable once the contract is fixed; differences show in wiring review, not build time \\
\hline
\end{tabularx}
\end{table}

\section{Use Case and Guardrail Architecture}

\subsection{Use case}
The use case is natural-language-driven firewall-rule provisioning: an operator requests, in natural language, that access be allowed from a source subnet to a destination subnet and port, and the pipeline must interpret, validate, check against policy, plan, simulate, obtain approval, apply, verify and, if verification fails, roll back the change. NetBox is the intended inventory source of truth for the overlap check at step 3; a mock firewall API, logged to a plain-text access log, stands in for a production device at steps 6, 8 and 9. Table~\ref{tab:usecase} gives the ten steps. Each step is implemented, and its guardrail independently verified, on all three platforms; only LangGraph and n8n connect all ten into one executable pipeline. Langflow's steps 1--6 remain independently verified sub-chains rather than a single connected run, for the scheduling reason given in Section V.

A local NetBox instance could not be brought up (a WSL2/Docker environment fault, unrelated to the platforms under test), so step 3 currently returns a fixed \texttt{overlap\_found: false} on all three platforms -- a stated implementation limitation, not a finding: its guardrail checks schema validity only, not real inventory data, in the results in Section VI.

We revised the step order used in an earlier draft of this pipeline after noting that a rollback plan is only useful if it exists \emph{before} the change is risked, not after. The planner (step 5) therefore emits the candidate rule and its rollback plan in the same structured output; step 10 executes that stored plan only if verification (step 9) fails.

\begin{table}[!t]
\caption{Guardrail-Checked Use-Case Steps}
\label{tab:usecase}
\centering
\footnotesize
\begin{tabularx}{\columnwidth}{@{}c X X@{}}
\hline
\textbf{\#} & \textbf{Action} & \textbf{Guardrail} \\
\hline
1  & Parse request into structured intent & Schema + field completeness \\
2  & Validate address/port/protocol syntax & Deterministic syntax check \\
3  & Query NetBox for overlapping rules & Schema only (fixed stub, Section IV.A) \\
4  & Policy/security check & CIDR width + port sensitivity ground truth \\
5  & Propose rule, change plan \emph{and} rollback plan & Schema + plan completeness \\
6  & Dry-run / simulate the change & Simulation result check \\
7  & Pause, request human approval & N/A (human decision) \\
8  & Apply the rule via the (mock) API & HTTP status check \\
9  & Verify the change took effect & Independent re-read of device state \\
10 & Execute rollback plan \emph{iff} step 9 failed & Conditional, deterministic \\
\hline
\end{tabularx}
\end{table}

\subsection{Guardrail architecture}
Every LLM-driven step is checked against a JSON Schema and, where applicable, against a ground truth computed independently of the model -- for example, CIDR prefix width and port membership in a sensitive-port set are computed in Python with the standard \texttt{ipaddress} library, never taken on the model's word. We call this the \emph{structural} layer.

Two of the four LLM-driven steps go further. Validation (step 2) and the policy check (step 4) ask the model not just to extract or generate content, but to render a yes/no verdict -- \texttt{is\_valid} and \texttt{policy\_pass}, respectively. For these two steps, the schema's own verdict field is re-checked as a second, \emph{business-decision} layer, independent of the structural layer: an output can be well-formed and internally consistent while still carrying a negative verdict, and the pipeline must stop on that verdict regardless. Intent parsing (step 1) and planning (step 5) carry no such field -- they extract or generate content, they do not adjudicate -- so a structural pass is the only gate they require by design, not by omission.

Where a business-decision field exists, both layers must pass for the pipeline to proceed; either one failing routes the run to a terminal rejection state that names the failing step. Section VII.A shows that this two-layer design was not wired identically across platforms: LangGraph enforces both layers at both business-decision steps, while n8n's branching nodes at both of these steps read only the structural layer.

\section{Experimental Methodology}

The LangGraph/n8n results reported in Section VI were collected with a single, fixed LLM -- \texttt{qwen2.5:7b}, served locally via Ollama -- so that any behavioral difference between platforms is attributable to the orchestration layer or to model stochasticity, not to a different model. This isolates the two failure modes discussed in Section VII without a second variable. To test whether those failure modes are specific to \texttt{qwen2.5:7b} or generalize across model families, we additionally repeated the LangGraph harness against GLM-5.1, a reasoning model hosted by the Poznan Supercomputing and Networking Center (PSNC); results are in Section VI.A.

Four intents, identical across platforms and held fixed throughout, probe distinct guardrail paths:
\begin{itemize}
\item \emph{simple} -- a syntactically complete, policy-compliant request (expected: approved);
\item \emph{edge\_incomplete} -- a request missing required fields (expected: rejected at step 1);
\item \emph{policy\_violating} -- a request with an overly wide source subnet and a sensitive destination port (expected: rejected at step 4);
\item \emph{guaranteed\_safe} -- a request designed to satisfy neither of the two defined high-risk conditions (expected: approved).
\end{itemize}

Each intent was run ten times per platform, with the human-approval step auto-approved to isolate LLM/guardrail stability from the (deterministic) approval decision itself, yielding 80 LangGraph/n8n runs. Full end-to-end wiring of steps 1--6 proved impractical to complete for Langflow within the project schedule: each step exists as an independently verified sub-chain, but connecting them would have required rebuilding prompt payloads and re-wrapped guardrail queries across roughly six non-trivial integration points. We therefore scoped the Langflow trial to its connected segment, steps 7--10 (human approval through conditional rollback), run ten times against a fixed static approval request. This segment contains no LLM call, so it measures infrastructure stability (suspend/resume reliability) rather than model stochasticity; this is a stated limitation, not a platform capability judgment.

Outcome detection relies on ground truth external to each platform's own reporting: for Langflow, the \texttt{outputs} dictionary returned by the v2 API was found to omit vertex keys inconsistently across otherwise-identical runs, so outcomes are instead read from the mock firewall's access log (presence/absence of \texttt{POST /rules/apply}, \texttt{GET /rules/\{id\}}, \texttt{DELETE /rules/\{id\}}).

\section{Results}

Table~\ref{tab:results} reports outcomes for the four intents on LangGraph and n8n (n=10 each); Langflow results are reported separately in Section VII.C.

\begin{table}[!t]
\caption{Outcomes per Intent, LangGraph vs.\ n8n (n=10 each)}
\label{tab:results}
\centering
\footnotesize
\begin{tabularx}{\columnwidth}{@{}X c c@{}}
\hline
\textbf{Intent (expected outcome)} & \textbf{LangGraph} & \textbf{n8n} \\
\hline
\texttt{edge\_incomplete} (reject, step 1) & 10/10 correct & 10/10 correct \\
\texttt{policy\_violating} (reject, step 4) & 10/10 correct & \textbf{0/10 correct} \\
\texttt{simple} (approve) & 0/10 correct & 2/10 correct \\
\texttt{guaranteed\_safe} (approve) & 0/10 correct & 3/10 correct \\
\hline
\end{tabularx}
\end{table}

Structural rejection (\texttt{edge\_incomplete}) is perfectly consistent across platforms: both reject the incomplete request in every trial, in under 6 s on average, before any downstream LLM call is made. This is the cheapest and most reliable guardrail in the pipeline.

Semantic rejection (\texttt{policy\_violating}) is not: LangGraph rejects it in all 10 trials; n8n approves it in 8 of 10, with the remaining 2 failing later, at the planner step, for unrelated reasons. Section VII.A traces this to an architectural cause rather than a model difference.

Approval of safe requests is unreliable on both platforms. \texttt{guaranteed\_safe} was designed to satisfy neither high-risk condition in the policy prompt, yet LangGraph approved it in 0 of 10 trials (10/10 rejected at the policy check) and n8n in only 3 of 10 (5 rejected at the policy check, 1 at the planner, 1 a malformed-JSON error rather than a guardrail rejection). A controlled ablation -- rerunning \texttt{guaranteed\_safe} on LangGraph with an explicit prompt instruction to default to approval absent a defined high-risk condition -- did not close the gap (0/10 approved across a further 10 trials), indicating the false-rejection rate is not attributable to a single fixable prompt defect. Section VII.B discusses this.

\subsection{Multi-model round (LangGraph)}

To test whether the two failure modes in Sections VII.A and VII.B are specific to \texttt{qwen2.5:7b} or generalize across model families, we repeated the same 40-trial LangGraph harness against GLM-5.1, a PCSS-hosted reasoning model (Section V). Table~\ref{tab:multimodel} reports both models' outcomes.

\begin{table}[!t]
\caption{LangGraph Outcomes per Intent, qwen2.5:7b vs.\ GLM-5.1 (n=10 each)}
\label{tab:multimodel}
\centering
\footnotesize
\begin{tabularx}{\columnwidth}{@{}X c c@{}}
\hline
\textbf{Intent} & \textbf{qwen2.5:7b} & \textbf{GLM-5.1} \\
\hline
\texttt{edge\_incomplete} & 10/10 correct & 10/10 correct \\
\texttt{policy\_violating} & 10/10 correct & 10/10 correct \\
\texttt{simple} & 0/10 correct & \textbf{10/10 correct} \\
\texttt{guaranteed\_safe} & 0/10 correct & \textbf{10/10 correct} \\
\hline
\end{tabularx}
\end{table}

GLM-5.1 reaches a perfect 40/40 correct outcomes, against qwen2.5:7b's 20/40 -- identical on the two structural/semantic-rejection intents, but with none of the false-rejection pattern on deliberately safe requests. Section VII.B discusses this directly.

\section{Discussion}

\subsection{A platform-level guardrail gap, not a model difference}
For \texttt{policy\_violating}, the model behaves identically on both platforms, correctly emitting \texttt{policy\_pass: false} with \texttt{risk\_level: "high"}. The divergence is downstream: LangGraph's guardrail node checks both layers from Section IV.B -- structural consistency and the model's own \texttt{policy\_pass} field -- as ordinary attribute reads on one state object. n8n's branching (If) node checks only the structural layer, since the guardrail is an HTTP call reporting structural consistency, and \texttt{policy\_pass} was never independently routed to the branching node. A well-formed, correctly-labeled-dangerous output is consequently read as ``passed'' and the run proceeds to apply.

This is not an isolated slip. n8n's validator node shows the identical pattern, reading only the structural \texttt{passed} field, never \texttt{is\_valid}: \texttt{rejected\_at\_validator} never appears across 40 n8n trials (Table~\ref{tab:results}), while LangGraph's validator check fired in 2 of 10 \texttt{policy\_violating} trials, catching what the policy check would otherwise have to catch downstream. Both n8n gates were wired the same way, so both fail the same way -- a systematic property of how n8n routes an LLM's own verdict, not a one-off mistake, and exactly the error the two-layer discipline in Section IV.B was designed to prevent. Its absence in one of two full implementations is itself evidence for this paper's central claim: the orchestration layer, not only the model, determines whether a guardrail design survives implementation.

\subsection{Model stochasticity as a distinct failure mode}
Independent of the platform gap, both platforms show a high false-rejection rate on requests engineered to be safe. Guardrail logs attribute this to two mechanisms: the model classifying a request as high-risk against satisfied-by-construction criteria (a business-decision error), and emitting an internally inconsistent output, e.g.\ \texttt{risk\_level: "high"} paired with \texttt{policy\_pass: true} (a structural error caught by the consistency check). An explicit prompt correction did not reduce the rejection rate, suggesting this is a property of the 7B-parameter model's reliability on structured risk classification, not of prompt phrasing.

This hypothesis holds. GLM-5.1, a larger PCSS-hosted reasoning model, reaches 40/40 correct outcomes on the identical harness (Table~\ref{tab:multimodel}), with zero false rejections on either safe intent, while agreeing with qwen2.5:7b on both rejection intents. The over-cautiousness bias is therefore a property of the smaller local model, not of the guardrail architecture, the prompt, or the pipeline design -- a materially different, and more actionable, conclusion than either failure mode viewed in isolation.

\subsection{Langflow: infrastructure stability under a narrower scope}
The Langflow segment (steps 7--10, ten trials, static request, no LLM call) completed successfully in all 10 trials, with a mean duration of 146 s per run (range: 109--240 s). The suspend/resume mechanism -- job submission, polling to a \texttt{suspended} state, retrieval of a pending request id, resume with an explicit decision, polling to completion -- proved reliable once the correct v2 endpoint was identified; an earlier attempt using the platform's synchronous build endpoint silently bypassed the pause entirely and is documented as a methodological hazard for any future use of this API. Because this segment carries no LLM call, its stability is not directly comparable to the LangGraph/n8n guardrail results in Table~\ref{tab:results}; it demonstrates that the underlying human-in-the-loop infrastructure is dependable, without speaking to end-to-end guardrail behavior across the full pipeline.

\subsection{Implications}
All three platforms support a supervisor pattern, where one agent delegates to specialized workers; we do not use it, and did not test it, so this is a design rationale, not a result. A guardrail attributes a failure to one named step only because exactly one agent call produced the tested output; a supervisor reasoning across steps before emitting output would blur that attribution. We judge this per-step auditability worth four LLM calls per trial rather than fewer.

Three practical guidelines follow. First, a guardrail specification is not self-enforcing: platforms that separate an agent's output from its own decision field across nodes need an explicit, auditable wire between them, checked at implementation review, not assumed from a guardrail call's mere presence. Second, false rejection of safe input is, on this evidence, at least as large an obstacle as missed detection of unsafe input, and deserves equal attention in guardrail design for smaller, locally-hosted models. Third, a supervisor pattern that merges steps should carry a correspondingly finer-grained guardrail on its internal reasoning, not just its final output.

\section{Conclusion and Future Work}
We implemented an identical ten-step, guardrail-checked firewall-provisioning pipeline in n8n, Langflow and LangGraph, and ran 130 trials across four intents and two models. The results separate two failure modes that are easy to conflate: LLM stochasticity, and a platform-specific n8n wiring gap that let a policy-violating request reach apply in 8 of 10 trials. Langflow's HITL/rollback infrastructure was stable across all 10 trials of its narrower, LLM-free scope. Critically, the stochasticity failure did not generalize: GLM-5.1 reached 40/40 correct outcomes on the identical LangGraph harness where \texttt{qwen2.5:7b} scored 20/40, showing the bias is a property of the smaller local model, not the guardrail architecture.

Remaining work includes a live NetBox instance in place of the stub (Section IV.A), completing Langflow's end-to-end wiring, testing the supervisor-pattern rationale (Section VII.D) empirically, and extending the multi-model round to n8n. Code, guardrail functions and the raw experiment log are public at \url{https://github.com/sandrailievskaa/telfor-agent-orchestration}.

\section*{Acknowledgment}
The additional open-source model evaluation described in Section V uses compute infrastructure provided by the Poznan Supercomputing and Networking Center (PSNC), Poland, made available through R. {\L}apacz. The authors thank PSNC for this access.

\begin{thebibliography}{00}

\bibitem{hong2025survey} J. Hong, N. Van Tu, and J. W.-K. Hong, ``A Comprehensive Survey on LLM-Based Network Management and Operations,'' \textit{International Journal of Network Management}, vol. 35, no. 6, e70029, 2025.

\bibitem{boateng2025survey} G. O. Boateng, H. Sami, A. Alagha, H. Elmekki, A. Hammoud, R. Mizouni, A. Mourad, H. Otrok, J. Bentahar, S. Muhaidat, C. Talhi, Z. Dziong, and M. Guizani, ``A Survey on Large Language Models for Communication, Network, and Service Management: Application Insights, Challenges, and Future Directions,'' \textit{IEEE Communications Surveys \& Tutorials}, 2025.

\bibitem{futureinternet2026} ``LLM-Based Multi-Agent Orchestration: A Survey of Frameworks, Communication Protocols, and Emerging Patterns,'' \textit{Future Internet}, vol. 18, no. 6, art. 326, 2026.

\bibitem{wang2025confucius} Z. Wang, S. Lin, G. Yan, S. Ghorbani, M. Yu, J. Zhou, N. Hu, L. Baruah, S. Peters, S. Kamath, J. Yang, and Y. Zhang, ``Intent-Driven Network Management with Multi-Agent LLMs: The Confucius Framework,'' in \textit{Proc. ACM SIGCOMM 2025 Conference}, pp. 347--362.

\bibitem{brodimas2025intent} D. Brodimas, A. Birbas, D. Kapolos, and S. Denazis, ``Intent-Based Infrastructure and Service Orchestration Using Agentic-AI,'' \textit{IEEE Open Journal of the Communications Society}, vol. 6, pp. 7150--7168, 2025.

\bibitem{zaid2025multiagent} H. Zaid, P. Safari, B. Shariati, A. Jafari, M. Balanici, and J. K. Fischer, ``Multi-Agent Design for LLM-assisted Network Management,'' in \textit{Proc. Optical Fiber Communication Conference (OFC) 2025}, paper W1A.4.

\end{thebibliography}

\end{document}
